\section{Chapter 5: Feedback control \texorpdfstring{\fr{162--218}}{}}
\label{ch5}

\subsection*{Big picture}
Chapter 4 told us how to \emph{recognise} stability. Chapter 5 is about \emph{creating} it by feedback, using only algebra on transfer functions:
\begin{enumerate}
\item test polynomials for stability without finding roots (Maxwell's problem);
\item the ring $\mathcal S$ of stable rational functions;
\item Nyquist's criterion (stability under feedback from a picture);
\item the simple feedback loop $T=\frac{G}{1+GK}$, PID controllers;
\item a guaranteed recipe to stabilise \emph{any} rational plant: coprime factorisation $\to$ Youla controller $K=X/Y$;
\item matrix tools: positive definite matrices, Lyapunov's criterion, Sylvester's equation.
\end{enumerate}

\subsection{Stable polynomials \texorpdfstring{\fr{163--164}}{}}\label{sec:stabpoly}\label{ss:5-1}
\begin{plainbox}
A polynomial is \textbf{stable} if all its roots are in the left half plane (so the matching terms $e^{\lambda t}$ decay).
Quick test without finding roots: a stable polynomial must have \textbf{all coefficients positive}. That is enough for degree 2,
but not for degree 3 or more (use $a_2a_1>a_0$ for degree 3).
\par\smallskip\textbf{In the exam this looks like:} ``Is this polynomial stable?'', ``for which $k$ is it stable?'', ``prove all coefficients are positive''.
\end{plainbox}
\begin{defbox}
A polynomial $h$ is \textbf{stable} if all its zeros lie in the open left half plane LHP $=\{s:\Re s<0\}$.
(So $h$ stable $\iff$ $1/h$ is a stable rational function.)
\end{defbox}
\begin{thmbox}[title={Proposition 0.43 (necessary condition)}]
If $h$ is a \textbf{monic real} stable polynomial then \textbf{all its coefficients are positive}.
\end{thmbox}
\begin{proofbox}[title={Proof of Proposition 0.43}]
Real coefficients $\Rightarrow$ $h(\lambda)=0\iff h(\bar\lambda)=0$. So the roots are real $\mu_j<0$, or conjugate pairs $\lambda_k,\bar\lambda_k$ with $\Re\lambda_k<0$, giving
$(s-\lambda_k)(s-\bar\lambda_k)=s^2-2\Re\lambda_k\,s+|\lambda_k|^2$. Hence
$h(s)=\prod_j(s-\mu_j)\prod_k(s^2-2\Re\lambda_k\,s+|\lambda_k|^2)$ with $-\mu_j>0$, $-2\Re\lambda_k>0$, $|\lambda_k|^2>0$.
Multiplying out polynomials with positive coefficients gives positive coefficients.
\end{proofbox}
\begin{warnbox}
The converse is \textbf{false} for degree $\ge3$: $s^3+s^2+s+1=(s+1)(s^2+1)$ has positive coefficients but roots $\pm i$. The AE2 polynomial $s^4+2s^3+s^2+4s+2$ also has positive coefficients but is unstable.
So use positivity only to \emph{rule out} stability (a zero or negative coefficient $\Rightarrow$ not stable).
\end{warnbox}
\begin{thmbox}[title={Useful tests by degree (degree 3 is Routh's extension, \emph{not proved} in the notes)}]
\begin{itemize}
\item Degree 1: $s+a_0$ stable $\iff a_0>0$.
\item Degree 2: $s^2+a_1s+a_0$ stable $\iff a_1>0$ and $a_0>0$. (Proof: roots $\frac{-a_1\pm\sqrt{a_1^2-4a_0}}2$; this is the damped oscillator table.)
\item Degree 3: $s^3+a_2s^2+a_1s+a_0$ stable $\iff a_2,a_1,a_0>0$ and $\boxed{a_2a_1>a_0}$.
\item Any degree: form the companion matrix and compute its eigenvalues numerically \fr{163}.
\end{itemize}
If you use the degree-3 test, say ``by the Routh--Hurwitz criterion''; when $a_2a_1=a_0$ the polynomial factorises as $(s+a_2)(s^2+a_1)$ with roots on the imaginary axis.
\end{thmbox}

\subsection{The ring \texorpdfstring{$\mathcal S$}{S} of stable rational functions \texorpdfstring{\fr{165--169}}{}}\label{ss:5-2}
\begin{plainbox}
$\mathcal S$ is the set of ``safe'' transfer functions: \textbf{proper} (does not blow up at high frequency) and all poles on the
left (does not blow up in time). Adding or multiplying safe functions gives safe functions.
\par\smallskip\textbf{In the exam this looks like:} ``Is $f\in\mathcal S$?'' or ``show $\mathcal S$ is not a field''.
\end{plainbox}
\begin{defbox}
$f=g/h$ is a \textbf{stable rational function}, $f\in\mathcal S$, if (i) $f$ is \textbf{proper} ($\deg g\le\deg h$) and (ii) all poles lie in the open LHP ($h$ stable).
\end{defbox}
\begin{thmbox}[title={Proposition 0.44}]
$\mathcal S$ is a commutative differential ring (closed under $+$, $\times$, scalar multiples and $d/ds$, Leibniz rule; no zero divisors on RHP). But $\mathcal S$ is \textbf{not a field}:
$\frac1{s+1}\in\mathcal S$ but its inverse $s+1$ is not proper, so $s+1\notin\mathcal S$.
\end{thmbox}
\begin{proofbox}[title={Proof sketch (closure)}]
Product: $\frac{g_1g_2}{h_1h_2}$ has $\deg(g_1g_2)\le\deg(h_1h_2)$ and the zeros of $h_1h_2$ are zeros of $h_1$ or $h_2$, so in the LHP.
Sum: $\frac{g_1h_2+g_2h_1}{h_1h_2}$, same reasoning. Derivative: by partial fractions $f=q+\sum a_j(s-\lambda_j)^{-n_j}$ with $q$ constant and $\Re\lambda_j<0$, so
$f'=\sum -n_ja_j(s-\lambda_j)^{-n_j-1}$ is strictly proper with the same poles.
\end{proofbox}

\subsection{Nyquist's criterion \texorpdfstring{\fr{174--177}}{}}\label{ss:5-3}
\begin{plainbox}
Look at the picture of $R(i\omega)$. If it never loops around the point $-1$, then closing the loop (feeding the output back
negatively) keeps the system stable. Where the curve crosses the negative real axis tells you how much extra gain you can add
before it becomes unstable.
\par\smallskip\textbf{In the exam this looks like:} ``Use Nyquist's criterion to find the values of $k$ for which the closed loop is stable''.
\end{plainbox}
\begin{thmbox}[title={Nyquist's criterion}]
Let the plant $R$ be \textbf{stable}. If the Nyquist contour $R(i\omega)$, $-\infty\le\omega\le\infty$, does \textbf{not pass through or wind around $-1$}, then
$T=\dfrac{R}{1+R}$ is stable, i.e.\ the feedback loop with constant feedback $-1$ is stable.\\
\textbf{Argument principle:} $N=Z-P$ where $N$ = net number of \textbf{clockwise} windings of $R(i\omega)$ round $-1$, $Z$ = number of zeros of $1+R$ in the RHP, $P$ = number of poles of $1+R$ in the RHP.
Here $P=0$ ($R$ stable) and $N=0$, so $Z=0$: no closed-loop poles in the RHP.
\end{thmbox}
\begin{proofbox}[title={Why $T$ is proper and where its poles can be \fr{175}}]
Write $R=c+p/q$ with $\deg p<\deg q$, $c=\lim_{s\to\infty}R\neq-1$. Then $T=\dfrac{cq+p}{(1+c)q+p}$, and $\deg((1+c)q+p)=\deg q$, so $T$ is proper.
A pole of $R$ gives a finite value of $T=R/(1+R)$, so the only possible poles of $T$ are zeros of $1+R$; the argument principle locates them.
\end{proofbox}
\begin{methodbox}[title={Method: use Nyquist to find a gain margin}]
\begin{enumerate}
\item Find where $R(i\omega)$ crosses the negative real axis: solve $\Im R(i\omega)=0$ (equivalently $\phi(\omega)=-\pi$).
\item Compute $R(i\omega)$ there. If it lies to the right of $-1$, no winding $\Rightarrow$ stable feedback loop.
\item With a gain $k$, the crossing point scales by $k$: stable while $k\,|R(i\omega_c)|<1$.
\end{enumerate}
\end{methodbox}
\begin{workbox}[title={Worked example: $R(s)=\dfrac{k}{(s+1)^3}$, $k>0$}]
\textbf{Nyquist.} $\phi(\omega)=-3\tan^{-1}\omega=-\pi$ at $\omega=\sqrt3$, where $|R|=\frac{k}{(1+3)^{3/2}}=\frac k8$. The plot crosses the negative real axis at $-k/8$, so it avoids winding round $-1$ iff $k<8$.\\
\textbf{Check algebraically.} Closed-loop denominator $(s+1)^3+k=s^3+3s^2+3s+(1+k)$: Routh needs $3\cdot3>1+k$, i.e.\ $k<8$ \checkmark. At $k=8$: $(s+3)(s^2+3)$, roots $\pm i\sqrt3$ (exactly the crossing frequency).
\end{workbox}

\subsection{The simple feedback loop and PID control \texorpdfstring{\fr{179--185}}{}}\label{sec:pid}\label{ss:5-4}
\begin{plainbox}
A \textbf{controller} watches the output and corrects the input, like cruise control. With negative feedback through $K$, the new
transfer function is $T=G/(1+GK)$. We choose $K$ (often $a+b/s+cs$, a PID) so that the new denominator is a stable polynomial.
\par\smallskip\textbf{In the exam this looks like:} ``Find a controller of the form $-a-b/s$ that stabilises $G$'' --- set up the denominator and choose $a,b$.
\end{plainbox}
\begin{center}
\begin{tikzpicture}
\node (u) at (0,0) {$u$}; \node[sumj] (s) at (1.3,0) {$+$}; \node[blk] (g) at (3,0) {$G$}; \node[tap] (t) at (4.4,0) {}; \node (y) at (5.4,0) {$y$};
\node[blk] (k) at (3,1.2) {$-K$};
\draw[ar] (u)--(s); \draw[ar] (s)--(g); \draw[ar] (g)--(t); \draw[ar] (t)--(y); \draw[ar] (t)|-(k); \draw[ar] (k)-|node[left,pos=.8]{$-Ky$}(s);
\end{tikzpicture}
\end{center}
\begin{thmbox}[title={Simple feedback loop (SFL)}]
Plant $G$ ($m\times k$), controller $-K$ ($k\times m$, negative feedback). From $y=G(u-Ky)$:
\[\boxed{T=(I+GK)^{-1}G}\qquad\text{(SISO: }T=\frac{G}{1+GK}\text{).}\]
Poles of $G$: \textbf{open-loop poles}; poles of $T$: \textbf{closed-loop poles}. If $G=p/q$ and $K=k$ (constant), $T=\dfrac{p}{kp+q}$.
\textbf{PID controller:} $K(s)=a+\dfrac bs+cs$ (proportional + integral + derivative). PD: abrupt correction (Whitehead's torpedoes); PI: gentle.
\end{thmbox}
\begin{methodbox}[title={Method: choose a controller to stabilise $G$}]
\begin{enumerate}
\item Write $T=\dfrac{G}{1+GK}=\dfrac{1}{1/G+K}$ and clear fractions to get $T=\dfrac{\text{num}}{\text{closed-loop polynomial}}$.
\item Write the closed-loop polynomial with the controller parameters as unknowns.
\item Impose the stability test (positive coefficients for degree 2; Routh for degree 3).
\item Choose nice values (e.g.\ make the polynomial $(s+1)^2$), and check that $T$ is \textbf{proper}.
\end{enumerate}
\end{methodbox}
\begin{workbox}[title={Worked example \fr{185}: stabilise $G(s)=\dfrac{s}{(s-1)(s+3)}$ by a PI controller}]
Open-loop poles $1,-3$: unstable. Take $K(s)=-a-\frac bs$ in $T=\frac{G}{1+GK}$ (this is the notes' sign convention). Then
\[\frac{1}{G}+K=\frac{(s-1)(s+3)}{s}-a-\frac bs=\frac{s^2+2s-3-as-b}{s},\qquad T=\frac{s}{s^2+(2-a)s-(b+3)}.\]
Need $2-a>0$ and $-(b+3)>0$. Choose $a=0$, $b=-4$: $T=\dfrac{s}{s^2+2s+1}=\dfrac{s}{(s+1)^2}\in\mathcal S$. So $K(s)=\frac4s$ (pure integral feedback) works.
\end{workbox}
\begin{warnbox}
Sign conventions differ: the diagram has the block $-K$ and the formula uses $K$. State your convention once, e.g.\ ``with $T=G/(1+GK)$'', and stick to it.
\end{warnbox}
\begin{workbox}[title={Worked example \fr{184}: proportional feedback can fix an improper plant}]
$G=\frac{s^2}{bs+c}$ ($b,c>0$) is not proper. With $K=k>0$: $T=\dfrac{s^2}{ks^2+bs+c}$, proper, poles $\frac{-b\pm\sqrt{b^2-4ck}}{2k}$ in the LHP (positive coefficients, degree 2). Stable.
\end{workbox}

\subsection{Euclid, B\'ezout and matrix completion \texorpdfstring{\fr{186--193}}{}}\label{ss:5-5}
\begin{plainbox}
The Euclidean algorithm finds the \textbf{highest common factor} of two polynomials by repeated long division, exactly like for
whole numbers. Working backwards writes $1=PM+QN$ whenever $P,Q$ have no common root. This identity is the engine of the next subchapter.
\par\smallskip\textbf{In the exam this looks like:} ``Find polynomials $M,N$ with $PM+QN=1$'' or ``complete this matrix to an invertible one''.
\end{plainbox}
\begin{methodbox}[title={Euclidean algorithm and B\'ezout for polynomials}]
\begin{enumerate}
\item Divide: $P_0=Q_1P_1+P_2$ with $\deg P_2<\deg P_1$; then $P_1=Q_2P_2+P_3$, and so on until a remainder is $0$. The last nonzero remainder $P_n$ is the highest common factor.
\item \textbf{Work back}: $P_n=P_{n-2}-Q_{n-1}P_{n-1}$, substitute for $P_{n-1}$, etc., to get $P_n=MP_0+NP_1$.
\item If $P_n$ is a nonzero constant $r$, divide by $r$ to get $MP_0/r+NP_1/r=1$.
\end{enumerate}
MATLAB: \texttt{[g,M,N]=gcd(P,Q)}. All in exact arithmetic, no roots needed.
\end{methodbox}
\begin{thmbox}[title={Lemma 0.47 and Proposition 0.48}]
Nonzero $P,Q\in\CC[s]$ have \textbf{no common zeros} $\iff$ $\exists$ polynomials $M,N$ with $PM+QN=1$.\\
For $f,g\in\CC[s]$ TFAE: (1) $\begin{bmatrix}f&g\\h&k\end{bmatrix}$ can be completed to a matrix with polynomial inverse; (2) $fk-gh=\Delta\in\CC\setminus\{0\}$; (3) hcf$(f,g)=1$; (4) $f,g$ have no common zeros.
\end{thmbox}
\begin{proofbox}[title={Proof of Lemma 0.47}]
($\Leftarrow$) If $\lambda$ were a common zero, $1=P(\lambda)M(\lambda)+Q(\lambda)N(\lambda)=0$, contradiction.
($\Rightarrow$) Euclid gives $PM+QN=R$ with $R$ the hcf. If $R=r$ is a nonzero constant, divide by $r$. Otherwise $\deg R\ge1$, so $R$ has a zero $\lambda$ (fundamental theorem of algebra); $R$ divides $P$ and $Q$, so $\lambda$ is a common zero, contradiction.
\end{proofbox}
\begin{workbox}[title={Worked examples}]
\textbf{(a) hcf.} $P_0=s^3+2s^2+s+2$, $P_1=s^2+3s+2$: $P_0=(s-1)P_1+(2s+4)$; $P_1=(\frac s2+\frac12)(2s+4)+0$. hcf $=2s+4\sim s+2$, and $2s+4=P_0-(s-1)P_1$.\\
\textbf{(b) B\'ezout.} $P=s+2$, $Q=s^2+1$: $Q=(s-2)P+5$, so $5=Q-(s-2)P$ and $\boxed{1=\tfrac15Q-\tfrac{s-2}5P}$. (Check: $\frac{s^2+1-(s^2-4)}5=1$ \checkmark)\\
\textbf{(c) Completion.} With $f=s+2$, $g=s^2+1$ take $k=\frac{2-s}5$, $h=-\frac15$: $fk-gh=\frac{4-s^2+s^2+1}5=1$, so $F=\begin{bmatrix}s+2&s^2+1\\-\frac15&\frac{2-s}5\end{bmatrix}$ has inverse $\begin{bmatrix}\frac{2-s}5&-(s^2+1)\\\frac15&s+2\end{bmatrix}$.\\
\textbf{(d) Integers} \fr{191}: $\begin{bmatrix}3&5\\c&d\end{bmatrix}$ has an integer inverse iff $3d-5c=\pm1$, e.g.\ $c=1,d=2$; possible because $\gcd(3,5)=1$.
\end{workbox}

\subsection{Coprime factorisation in \texorpdfstring{$\mathcal S$}{S} \texorpdfstring{\fr{194--200}}{}}\label{sec:coprime}\label{ss:5-6}
\begin{plainbox}
Goal: a \textbf{recipe that stabilises any plant}. Write $G=P/Q$ using only ``safe'' functions (the substitution $\lambda=1/(1+s)$ does this),
find $X,Y$ with $PX+QY=1$ (Euclid), and take the controller $K=X/Y$. Then all closed-loop functions are products of safe functions, so safe.
\par\smallskip\textbf{In the exam this looks like:} ``Find a coprime factorisation and a stabilising controller, and verify''. A long question, but purely mechanical.
\end{plainbox}
\begin{defbox}
$P,Q\in\mathcal S$ are \textbf{coprime} if there are $X,Y\in\mathcal S$ with $PX+QY=1$.
\textbf{Proposition 0.49:} every rational $G$ can be written $G=P/Q$ with $P,Q$ coprime in $\mathcal S$, found in exact arithmetic.
\end{defbox}
\begin{methodbox}[title={The coprime factorisation algorithm (4 steps)}]
\begin{enumerate}
\item Write $G=M(s)/N(s)$ with polynomials $M,N$ having \textbf{no common zeros}; let $m=\max(\deg M,\deg N)$.
\item Substitute $s=\dfrac{1-\lambda}{\lambda}$ (i.e.\ $\lambda=\dfrac1{1+s}$) and clear denominators: $\tilde M(\lambda)=\lambda^mM\big(\frac{1-\lambda}\lambda\big)$, $\tilde N(\lambda)=\lambda^mN\big(\frac{1-\lambda}\lambda\big)$.
\item Euclid/B\'ezout in $\lambda$: find polynomials $\tilde X,\tilde Y$ with $\tilde M\tilde X+\tilde N\tilde Y=1$.
\item Put back $\lambda=\dfrac1{1+s}$: $P=\tilde M,\ Q=\tilde N,\ X=\tilde X,\ Y=\tilde Y$ evaluated at $\frac1{1+s}$. All four are polynomials in $\frac1{1+s}$, so proper with only pole $-1$: they lie in $\mathcal S$.
\end{enumerate}
\end{methodbox}
\begin{thmbox}[title={Proposition 0.50 and Theorem 0.52 (Youla)}]
If $G=P/Q$ with $PX+QY=1$ in $\mathcal S$, then the controller $\boxed{K=X/Y}$ gives
\[\frac{1}{1+GK}=QY,\quad\frac{G}{1+GK}=PY,\quad\frac{K}{1+GK}=XQ,\quad\frac{GK}{1+GK}=PX,\]
all in $\mathcal S$. So $K=X/Y$ \textbf{internally stabilises} the SFL. Every rational SISO can be stabilised by a rational (possibly unstable) controller.
\end{thmbox}
\begin{proofbox}[title={Proof (one line each)}]
$1+GK=1+\frac{PX}{QY}=\frac{QY+PX}{QY}=\frac1{QY}$. Hence $\frac1{1+GK}=QY$, $\frac{G}{1+GK}=\frac PQ\,QY=PY$, $\frac{K}{1+GK}=\frac XY\,QY=XQ$, $\frac{GK}{1+GK}=\frac{PX}{QY}QY=PX$.
Products of elements of $\mathcal S$ are in $\mathcal S$.
\end{proofbox}

\begin{workbox}[title={Worked example A: $G(s)=\dfrac1{s-1}$}]
\textbf{1.} $M=1$, $N=s-1$, $m=1$.\quad
\textbf{2.} $\tilde M=\lambda$, $\tilde N=\lambda\big(\frac{1-\lambda}\lambda-1\big)=1-2\lambda$.\quad
\textbf{3.} By inspection $2\tilde M+1\cdot\tilde N=1$: $\tilde X=2$, $\tilde Y=1$.\\
\textbf{4.} $P=\frac1{1+s}$, $Q=1-\frac2{1+s}=\frac{s-1}{s+1}$, $X=2$, $Y=1$. Check $P/Q=\frac1{s-1}$ \checkmark. Controller $K=X/Y=2$:
$\frac{G}{1+GK}=\frac1{s+1}$, $\frac1{1+GK}=\frac{s-1}{s+1}$, $\frac{K}{1+GK}=\frac{2(s-1)}{s+1}$, $\frac{GK}{1+GK}=\frac2{s+1}$, all in $\mathcal S$.
\end{workbox}

\begin{workbox}[title={Worked example B: $G(s)=\dfrac1{s^2}$ (double integrator)}]
\textbf{1.} $M=1$, $N=s^2$, $m=2$.\quad
\textbf{2.} $\tilde M=\lambda^2$, $\tilde N=(1-\lambda)^2=\lambda^2-2\lambda+1$.\\
\textbf{3.} Euclid: $\tilde N=1\cdot\tilde M+(1-2\lambda)$; \ $\tilde M=\lambda^2=(1-2\lambda)\big(-\frac\lambda2-\frac14\big)+\frac14$. Work back:
$\frac14=\tilde M+(1-2\lambda)\big(\frac\lambda2+\frac14\big)=\tilde M+(\tilde N-\tilde M)\big(\frac\lambda2+\frac14\big)=\tilde M\big(\tfrac34-\tfrac\lambda2\big)+\tilde N\big(\tfrac\lambda2+\tfrac14\big)$.
Times 4: $\tilde M(3-2\lambda)+\tilde N(1+2\lambda)=1$.\\
\textbf{4.} $P=\frac1{(1+s)^2}$, $Q=\frac{s^2}{(1+s)^2}$, $X=3-\frac2{1+s}=\frac{3s+1}{s+1}$, $Y=1+\frac2{1+s}=\frac{s+3}{s+1}$; controller $K=\dfrac{3s+1}{s+3}$.\\
\textbf{Check:} $1+GK=\frac{s^2(s+3)+3s+1}{s^2(s+3)}=\frac{(s+1)^3}{s^2(s+3)}$, so $\frac{G}{1+GK}=\frac{s+3}{(s+1)^3}$, $\frac1{1+GK}=\frac{s^2(s+3)}{(s+1)^3}$, $\frac{K}{1+GK}=\frac{s^2(3s+1)}{(s+1)^3}$, $\frac{GK}{1+GK}=\frac{3s+1}{(s+1)^3}$: all proper with poles only at $-1$ \checkmark
\end{workbox}

\begin{workbox}[title={Worked example C \fr{196}: $G(s)=\dfrac{s^2+5}{s^2-s+1}$}]
With $s=\frac{1-\lambda}{\lambda}$, $m=2$: $\tilde M=6\lambda^2-2\lambda+1$, $\tilde N=3\lambda^2-3\lambda+1$, and Euclid gives
$\big(\tfrac{24\lambda-2}{7}\big)\tilde N-\big(\tfrac{12\lambda-9}{7}\big)\tilde M=1$. So $\tilde Y=\frac{24\lambda-2}7$, $\tilde X=-\frac{12\lambda-9}{7}$, and with $\lambda=\frac1{1+s}$:
$Y=\frac{22-2s}{7(1+s)}$, $X=\frac{9s-3}{7(1+s)}$, giving the stabilising controller $K=\dfrac XY=\dfrac{9s-3}{22-2s}$ (itself unstable, pole at $s=11$, which is allowed).
Indeed $\frac{G}{1+GK}=\frac{-2(s-11)(s^2+5)}{7(s+1)^3}\in\mathcal S$.
\end{workbox}

\subsection{Well-posedness and internal stability \texorpdfstring{\fr{201--206}}{}}\label{sec:internal}\label{ss:5-7}
\begin{plainbox}
It is not enough for the output to behave: \textbf{every signal inside the loop} must stay bounded too. If the controller cancels an
unstable part of the plant, the problem is hidden, not removed. So we check all four functions, not just $G/(1+GK)$.
\par\smallskip\textbf{In the exam this looks like:} ``Is the loop internally stable?'', ``when is the SFL well posed?''.
\end{plainbox}
\begin{thmbox}[title={Lemma 0.51}]
SFL with inputs $r,d,n$ and states $x=r-z$, $y=d+Kx$, $z=n+Gy$, i.e.\ $\begin{bmatrix}1&0&1\\-K&1&0\\0&-G&1\end{bmatrix}\begin{bmatrix}x\\y\\z\end{bmatrix}=\begin{bmatrix}r\\d\\n\end{bmatrix}$.
It is \textbf{well posed} (inputs determine states uniquely) iff $1+GK\neq0$, and then
\[\begin{bmatrix}x\\y\\z\end{bmatrix}=\frac1{1+GK}\begin{bmatrix}1&-G&-1\\K&1&-K\\GK&G&1\end{bmatrix}\begin{bmatrix}r\\d\\n\end{bmatrix}=\Phi\begin{bmatrix}r\\d\\n\end{bmatrix}.\]
\textbf{Internally stable}: every entry of $\Phi$ is in $\mathcal S$, i.e.\ $\frac1{1+GK},\frac G{1+GK},\frac{GK}{1+GK},\frac K{1+GK}\in\mathcal S$.
\end{thmbox}
\begin{exambox}[title={Why we need all four (a classic exam trap)}]
$G=\frac{s-1}{s+1}$, $K=\frac{1}{s-1}$ (the controller cancels the plant's RHP zero). Then $GK=\frac1{s+1}$ and
$\frac1{1+GK}=\frac{s+1}{s+2}$, $\frac{G}{1+GK}=\frac{s-1}{s+2}$, $\frac{GK}{1+GK}=\frac{1}{s+2}$ are all stable, \textbf{but}
$\frac{K}{1+GK}=\frac{s+1}{(s-1)(s+2)}$ has a pole at $1$. So the loop is \textbf{not internally stable}: an unstable pole--zero cancellation hides inside the loop, and the internal signal $y$ responds to the reference $r$ through the unstable $\frac{K}{1+GK}$, so it blows up even though the output $z$ looks fine.
\end{exambox}

\subsection{Self-adjoint and positive definite matrices \texorpdfstring{\fr{207--209}}{}}\label{sec:posdef}\label{ss:5-8}
\begin{plainbox}
A \textbf{positive definite} matrix is the matrix version of a positive number: $\langle Kz,z\rangle>0$ for every nonzero $z$.
Three equivalent tests: the definition, all eigenvalues positive, or all top-left determinants positive (easiest by hand).
\par\smallskip\textbf{In the exam this looks like:} ``Show that $S'KS$ / $\exp(A')K\exp(A)$ / $K+L$ is positive definite'' (AE3, A3.4).
\end{plainbox}
\begin{thmbox}
\begin{itemize}
\item \textbf{Prop 0.53:} if $A=A'$ then the eigenvalues are real and $A=U\Delta U'$ with $U$ unitary ($U'U=I$) and $\Delta$ real diagonal; eigenvectors for distinct eigenvalues are orthogonal.
\item $K$ is \textbf{positive definite} if $K=K'$ and $\langle Kz,z\rangle>0$ for all $z\neq0$.
\item \textbf{Lemma 0.54:} for $K=K'$, TFAE: (1) $\langle Kz,z\rangle>0\ \forall z\neq0$; (2) all eigenvalues $\kappa_j>0$; (3) all \textbf{leading principal minors} $\Delta_j>0$.
\item The product of positive definite matrices is generally \textbf{not} positive definite.
\end{itemize}
\end{thmbox}
\begin{workbox}[title={Worked example (AE3, A3.4): all five parts}]
$K$ real symmetric positive definite.\\
\textbf{(i)} If $Kv=\kappa v$, $v\neq0$: $0<\langle Kv,v\rangle=\langle\kappa v,v\rangle=\kappa\|v\|^2$, so $\kappa>0$ (in particular real).\\
\textbf{(ii)} $K$ is diagonalisable (Prop 0.53), so $\det K=\prod\kappa_j>0$ and $\tr K=\sum\kappa_j>0$.\\
\textbf{(iii)} $(S'KS)'=S'K'S''=S'KS$, so self-adjoint. For $z\neq0$, $Sz\neq0$ ($S$ invertible), so $\langle S'KSz,z\rangle=\langle KSz,Sz\rangle>0$.\\
\textbf{(iv)} $\exp(A)'=\exp(A')$ (take adjoints term by term: $(A^k)'=(A')^k$), and $\exp(A)$ is invertible with inverse $\exp(-A)$. Apply (iii) with $S=\exp(A)$: $\exp(A')K\exp(A)$ is positive definite.\\
\textbf{(v)} $(K+L)'=K+L$ and $\langle(K+L)z,z\rangle=\langle Kz,z\rangle+\langle Lz,z\rangle>0$ for $z\neq0$.
\end{workbox}
\begin{workbox}[title={Worked example: leading principal minors}]
$K=\begin{bmatrix}2&-1&0\\-1&2&-1\\0&-1&2\end{bmatrix}$: $\Delta_1=2$, $\Delta_2=4-1=3$, $\Delta_3=2(4-1)-(-1)(-2)=4$. All $>0$, so positive definite.
\end{workbox}

\subsection{Lyapunov's criterion and Lyapunov's equation \texorpdfstring{\fr{210--218}}{}}\label{ss:5-9}
\begin{plainbox}
Lyapunov's idea is an ``energy'' $V=\langle KX,X\rangle$. If you can show the energy never increases, the state can never run away.
This checks stability \emph{without} finding eigenvalues.
\par\smallskip\textbf{In the exam this looks like:} ``Use $K=I$ to show solutions are bounded'' or ``prove Lyapunov's criterion''.
\end{plainbox}
\begin{thmbox}[title={Theorem 0.55 (Lyapunov's criterion)}]
If there is a positive definite $K$ such that $Q=-(A'K+KA)$ is also positive definite, then all solutions of $\dot X=AX$ are bounded on $(0,\infty)$.
\end{thmbox}
\begin{proofbox}[title={Proof of Lyapunov's criterion (learn this)}]
Let $V(t)=\langle KX(t),X(t)\rangle\ge0$. Using $\dot X=AX$:
\begin{align*}\frac{dV}{dt}&=\langle K\dot X,X\rangle+\langle KX,\dot X\rangle=\langle KAX,X\rangle+\langle KX,AX\rangle\\&=\langle KAX,X\rangle+\langle A'KX,X\rangle=-\langle QX,X\rangle\le0.\end{align*}
So $V$ is decreasing. If $\kappa_n>0$ is the smallest eigenvalue of $K$, then
$\kappa_n\|X(t)\|^2\le\langle KX(t),X(t)\rangle\le\langle KX_0,X_0\rangle$, so $\|X(t)\|\le\big(\langle KX_0,X_0\rangle/\kappa_n\big)^{1/2}$ for all $t\ge0$.
\end{proofbox}
\begin{workbox}[title={Worked example: try $K=I$}]
$A=\begin{bmatrix}-1&2\\0&-3\end{bmatrix}$, $K=I$: $Q=-(A^\top+A)=-\begin{bmatrix}-2&2\\2&-6\end{bmatrix}=\begin{bmatrix}2&-2\\-2&6\end{bmatrix}$; minors $2>0$, $12-4=8>0$. So $Q$ is positive definite and all solutions are bounded.
(If $K=I$ fails, try another $K$, or choose $Q$ and solve for $K$.)
\end{workbox}
\begin{thmbox}[title={Corollary 0.59 (Lyapunov's equation)}]
If all eigenvalues of $A$ are in the open LHP, then for every positive definite $P$,
$L=\displaystyle\int_0^\infty\exp(tA)P\exp(tA')\,dt$ is positive definite and solves $\boxed{AL+LA'=-P}$. (MATLAB: \texttt{L=lyap(A,P)}.)
\end{thmbox}
\begin{workbox}[title={Worked example: solve $AL+LA^\top=-I$ for $A=\begin{bmatrix}0&1\\-2&-3\end{bmatrix}$}]
Let $L=\begin{bmatrix}l_1&l_2\\l_2&l_3\end{bmatrix}$. $AL=\begin{bmatrix}l_2&l_3\\-2l_1-3l_2&-2l_2-3l_3\end{bmatrix}$ and $LA^\top=(AL)^\top$. Entries of $AL+(AL)^\top=-I$:
$(1,1)$: $2l_2=-1\Rightarrow l_2=-\frac12$; \ $(2,2)$: $-4l_2-6l_3=-1\Rightarrow l_3=\frac12$; \ $(1,2)$: $l_3-2l_1-3l_2=0\Rightarrow l_1=1$.
$L=\begin{bmatrix}1&-\frac12\\-\frac12&\frac12\end{bmatrix}$, minors $1$ and $\frac14$: positive definite, as the Corollary predicts (eigenvalues $-1,-2$).
\end{workbox}

\subsection{Sylvester's equation \texorpdfstring{\fr{212--216}}{}}\label{sec:sylv}\label{ss:5-10}
\begin{plainbox}
$AY+YB=-C$ is just a (big) system of linear equations for the entries of $Y$. It has exactly one solution for every $C$ unless
$A$ and $-B$ share an eigenvalue (like $ay+yb=-c$ fails when $a=-b$).
\par\smallskip\textbf{In the exam this looks like:} ``When does Sylvester's equation have a unique solution?''.
\end{plainbox}
\begin{thmbox}
$AY+YB=-C$ is a linear system of $n^2$ equations in $n^2$ unknowns. Either a unique solution for all $C$, or infinitely many for some $C$, or none for some $C$ (Prop 0.56).
\textbf{Prop 0.57:} unique solution for every $C$ $\iff$ $A$ and $-B$ have \textbf{no common eigenvalue}.
\textbf{Prop 0.58:} if all eigenvalues of $A$ and $B$ are in the open LHP, $Y=\int_0^\infty\exp(tA)C\exp(tB)\,dt$ is the unique solution (MATLAB \texttt{lyap(A,B,C)}).
\end{thmbox}
\begin{proofbox}[title={Proof of Prop 0.58}]
$\|\exp(tA)\|\le M_1e^{-\delta_1t}$, $\|\exp(tB)\|\le M_2e^{-\delta_2t}$ so the integral converges. Then
$AY+YB=\int_0^\infty\frac{d}{dt}\big(\exp(tA)C\exp(tB)\big)dt=\big[\exp(tA)C\exp(tB)\big]_0^\infty=0-C=-C$.
So $Y\mapsto AY+YB$ is surjective on the $n^2$-dimensional space $M_{n\times n}(\CC)$, hence injective: the solution is unique.
\end{proofbox}
\begin{workbox}[title={Worked example}]
$1\times1$ intuition: $ay+yb=-c$ has a unique solution iff $a+b\neq0$, i.e.\ $a\neq-b$.\\
$A=\mathrm{diag}(1,2)$, $B=\mathrm{diag}(-1,3)$: eigenvalues of $-B$ are $1,-3$; $1$ is shared with $A$. Indeed $Y=\begin{bmatrix}1&0\\0&0\end{bmatrix}$ gives $AY+YB=Y-Y=0$, so solutions are not unique.
\end{workbox}

\subsection{Exam-style questions for Chapter 5}\label{ss:5-11}
\begin{answerbox}
\textbf{Polynomial stability:} quote the test you use (positive coefficients / Routh--Hurwitz) and check each condition.\\
\textbf{Controller design:} state $T=G/(1+GK)$, find the closed-loop denominator, state conditions, choose values, give the final $T$ and say why it is in $\mathcal S$.\\
\textbf{Coprime/Youla:} follow the 4 steps with headings, show $PX+QY=1$ explicitly, give $K=X/Y$, list the four closed-loop functions.
\end{answerbox}
\begin{qbox}[title={Q5.1 (stable polynomials)}]
(a) Is $s^3+2s^2+3s+5$ stable? (b) For which real $k$ is $s^3+2s^2+3s+k$ stable? (c) Show $s^4+s^3-2s^2+s+1$ is not stable.
\end{qbox}
\begin{solbox}
(a) Coefficients positive and $2\cdot3=6>5$, so stable (Routh--Hurwitz). (b) Need $k>0$ and $6>k$: $0<k<6$. At $k=6$: $(s+2)(s^2+3)$, roots $\pm i\sqrt3$.
(c) The coefficient $-2$ is negative, so by Prop 0.43 it is not stable.
\end{solbox}

\begin{qbox}[title={Q5.2 (proportional feedback)}]
$G(s)=\dfrac{1}{s(s+1)(s+2)}$, $K=k>0$. Find the closed-loop transfer function and the values of $k$ giving a stable closed loop.
\end{qbox}
\begin{solbox}
$T=\frac{p}{kp+q}=\frac{1}{s^3+3s^2+2s+k}$. Routh: $3\cdot2>k$, so stable iff $0<k<6$. At $k=6$, $(s+3)(s^2+2)$: poles $\pm i\sqrt2$ (marginal).
\end{solbox}

\begin{qbox}[title={Q5.3 (coprime factorisation and Youla)}]
Let $G(s)=\dfrac{1}{s-2}$. Find $P,Q,X,Y\in\mathcal S$ with $G=P/Q$, $PX+QY=1$, and a controller that internally stabilises the SFL.
\end{qbox}
\begin{solbox}
$m=1$: $\tilde M=\lambda$, $\tilde N=\lambda(\frac{1-\lambda}\lambda-2)=1-3\lambda$, and $3\tilde M+\tilde N=1$.
$P=\frac1{1+s}$, $Q=1-\frac3{1+s}=\frac{s-2}{s+1}$, $X=3$, $Y=1$; $K=3$.
Check: $\frac1{1+GK}=\frac{s-2}{s+1}$, $\frac{G}{1+GK}=\frac1{s+1}$, $\frac{K}{1+GK}=\frac{3(s-2)}{s+1}$, $\frac{GK}{1+GK}=\frac{3}{s+1}$: all in $\mathcal S$.
\end{solbox}

\begin{qbox}[title={Q5.4 (Lyapunov)}]
Show using $K=I$ that every solution of $\dot X=AX$, $A=\begin{bmatrix}-2&1\\-1&-1\end{bmatrix}$, is bounded.
\end{qbox}
\begin{solbox}
$A^\top+A=\begin{bmatrix}-4&0\\0&-2\end{bmatrix}$, so $Q=-(A^\top+A)=\mathrm{diag}(4,2)$, positive definite. By Theorem 0.55 solutions are bounded.
(Indeed $\chi_A=s^2+3s+3$, eigenvalues $\frac{-3\pm i\sqrt3}{2}$, in the LHP.)
\end{solbox}
\newpage
